\chapter{Exponential Decay of Correlations}\label{ch:correlations}

This chapter records the transfer-matrix expression for thermodynamic-limit
correlators of a normal MPS and the resulting exponential decay estimate.
The key mechanism is spectral: after removing the leading eigenvalue
contribution, the connected part is governed by subleading eigenvalues of the
transfer map.

When the MPS tensor generates a parent Hamiltonian
(Chapter~\ref{ch:parent_hamiltonian}), the correlator quantifies spatial
correlations in the ground state.  The exponential decay rate is set by the
spectral properties of the transfer map, i.e.\ by the modulus of its
subleading eigenvalues.

\section{Connected correlators from the transfer map}

\begin{definition}[One-point expectation]\label{def:one_point_expectation}
    \lean{MPSTensor.onePointExpectation}
    \leanok
    \uses{def:transfer_map, thm:qpf}
    Fix an MPS tensor $A$ and a right fixed point $\rho^R$ of its transfer map
    (whose existence, and uniqueness up to scaling under the hypotheses of
    Theorem~\ref{thm:qpf}, are guaranteed by the quantum Perron--Frobenius theorem).
    For a local observable $X \in \MN{D}$ define
    \[
        \langle X_0 \rangle := \tr(X\rho^R).
    \]
\end{definition}

\begin{definition}[Two-point expectation at distance $n$]\label{def:two_point_expectation}
    \lean{MPSTensor.twoPointExpectation}
    \leanok
    \uses{def:transfer_map, def:one_point_expectation}
    For observables $X,Y \in \MN{D}$ and $n\ge 0$,
    \[
        \langle X_0 Y_n \rangle
        := \tr\!\left(Y\,\E_A^n(X\rho^R)\right).
    \]
    Equivalently: insert $X$ at site $0$, propagate by $\E_A^n$, then insert
    $Y$ and take the trace.
\end{definition}

\begin{definition}[Connected correlator]\label{def:connected_correlator}
    \lean{MPSTensor.connectedCorrelator}
    \leanok
    \uses{def:two_point_expectation, def:one_point_expectation}
    The connected two-point function is
    \[
        C(X,Y;n) := \langle X_0Y_n\rangle - \langle X_0\rangle\langle Y_0\rangle.
    \]
\end{definition}

\section{Spectral expansion and exponential decay}

\begin{theorem}[Sum of exponentials]\label{thm:correlator_sum_exponentials}
    \lean{MPSTensor.connectedCorrelator_eq_sum}
    \leanok
    \uses{def:connected_correlator, thm:spectral_radius_le_one, thm:qpf}
    For a normal MPS, the connected correlator admits a spectral expansion of
    the form
    \[
        C(X,Y;n) = \sum_{j=1}^{D^2-1} c_j\,\lambda_j^n,
    \]
    where $\lambda_j$ are subleading transfer-matrix eigenvalues.
\end{theorem}

\begin{proof}
    The leading eigenvalue contribution ($\lambda_1=1$) equals
    $\langle X_0\rangle\langle Y_0\rangle$ and cancels in the connected part.
    Apply spectral decomposition to $\E_A^n$.
\end{proof}

\begin{theorem}[Exponential decay bound]\label{thm:correlator_exponential_bound}
    \lean{MPSTensor.connectedCorrelator_bound}
    \leanok
    \uses{thm:correlator_sum_exponentials}
    If $|\lambda_j| \le |\lambda_2| < 1$ for all subleading eigenvalues, then
    there exists $C_{XY}\ge 0$ such that
    \[
        |C(X,Y;n)| \le C_{XY}\,|\lambda_2|^n.
    \]
\end{theorem}

\begin{proof}
    Combine the sum-of-exponentials expansion with the triangle inequality.
\end{proof}

\begin{definition}[Correlation length]\label{def:correlation_length}
    \lean{MPSTensor.correlationLength}
    \leanok
    \uses{thm:correlator_exponential_bound}
    The correlation length associated with $\lambda_2$ is
    \[
        \xi := -\frac{1}{\log|\lambda_2|}.
    \]
\end{definition}

\begin{lemma}[Correlation length is positive]\label{lem:correlation_length_pos}
    \lean{MPSTensor.correlationLength_pos}
    \leanok
    \uses{def:correlation_length}
    If $|\lambda_2| < 1$ and $\lambda_2 \ne 0$, then $\xi > 0$.
\end{lemma}

\begin{proof}
    \leanok
    Since $0 < |\lambda_2| < 1$, we have $\log|\lambda_2| < 0$, so
    $\xi = -1/\log|\lambda_2| > 0$.
\end{proof}

\section{Relation to parent Hamiltonians}

\begin{remark}\label{rem:correlations_parent_hamiltonian}
    \uses{def:parent_hamiltonian, lem:parent_hamiltonian_annihilates,
          def:correlation_length, thm:correlator_exponential_bound}
    When the MPS tensor $A$ generates a parent Hamiltonian $H_N(A,L)$
    (Definition~\ref{def:parent_hamiltonian}) whose ground state is the
    matrix product vector (Lemma~\ref{lem:parent_hamiltonian_annihilates}),
    the exponential decay bound
    (Theorem~\ref{thm:correlator_exponential_bound}) shows that ground-state
    correlations decay with a finite correlation length $\xi$.
\end{remark}

\begin{remark}\label{rem:zero_correlation_length}
    \uses{thm:spectral_radius_le_one, def:correlation_length}
    The identity $\E^2=\E$ is equivalent to vanishing subleading spectrum;
    equivalently, $\xi=0$ and connected correlations vanish at nonzero
    separation.
\end{remark}
